% Preamble for Pandoc Generation
% The page style is the base; the book is layered on top.
%
% latex/page-style.sty       the page: base packages, geometry, running-head frame
% latex/theorem-frame.sty    the mdframed style and the \declaretheorem key the
%                            book's environments are framed with
% latex/latex-template.sty   the book's macros, colours and helpers
% latex/chapter-style.sty    chapters, \part banners, running heads
% latex/theorem-envs.sty     the environments filters/theorems.lua emits, in that
%                            frame, in the colours from the book's config

% ======================================================================================
% ONE DECISION PER SHARED PACKAGE
% ======================================================================================
% Sixteen packages are loaded by both the page style and the book. Most take no options
% and so cannot clash; these are the ones that do. Naming the options here, before
% anything pulls the packages in, means no load order can produce an option clash
% and there is a single place that says which settings win.

% xcolor: page-style loads it plain, latex-template wants dvipsnames. dvipsnames wins.
\PassOptionsToPackage{dvipsnames}{xcolor}
% hyperref: loaded exactly once, by templates/latex/template.tex, after
% latex/macros.tex (it reclaims \H, restored at the bottom of this file). These
% are the book's link settings, handed to that one load.
\PassOptionsToPackage{pdfusetitle,linktoc=page,allcolors=blue}{hyperref}
% mdframed: theorem-frame and latex-template both ask for framemethod=TikZ, so
% the second load is a no-op. Do NOT hand it the same option from here as well:
% mdframed reads the frame method by \input-ing a file, and it would do it twice.
% tcolorbox (most,breakable,skins,theorems) is loaded only by latex-template,
% and caption's skip=1em only by page-style, which is loaded first -- latex-template's
% later option-free load of caption is a subset and cannot clash.

% ======================================================================================
% THE BASE
% ======================================================================================
% Page geometry, running-head frame, paragraph style, tables and figures, and
% the packages they need. latex/page-style.sty is the one place the page is decided; it
% keeps the book's uniform 2.5cm margin, so the measure -- and every line break
% in 3,300 pages -- is unchanged.
\usepackage{latex/page-style}

\usepackage{bbm}

% ======================================================================================
% THE BOOK ON TOP
% ======================================================================================
% Macros, colours and helpers. page-style owns the page (skipgeometry), the pandoc
% template owns hyperref (skiphyperref), latex/chapter-style owns the headings
% (skiptitlesec) and latex/theorem-envs owns the environments (skipenv).
\usepackage[skipgeometry,skiphyperref,skiptitlesec,skipenv]{latex/latex-template}
\let\maketitle\relax

% List spacing (1.2x)
\usepackage{enumitem}
\setlist{itemsep=1.2em, parsep=0pt}

% Chapter styling (bookmode: two-level no box, three-level black box, correct font
% sizes) plus the \part banners and running heads, none of which page-style has.
\usepackage[bookmode]{latex/chapter-style}

% The environments filters/theorems.lua emits, in the book's framed style. Which
% environments exist, their names, counters and colours come from the book's
% config: the build writes them to _build/tmp/environments.tex, included right
% after this file. That covers the numbered forms, the unnumbered ones and the
% starred variants the preface uses (definition*, theorem*, ...).
\usepackage{latex/theorem-envs}

% Code listings (runnable cells) and function plots (components.lua)
\usepackage{listings}
\lstset{basicstyle=\ttfamily\small, columns=fullflexible, keepspaces=true, breaklines=true,
  frame=single, rulecolor=\color{black!25}, framesep=6pt, xleftmargin=6pt, xrightmargin=6pt,
  keywordstyle=\color{blue!60!black}\bfseries, commentstyle=\color{black!55}\itshape,
  stringstyle=\color{red!60!black}, showstringspaces=false, upquote=true}
\usepackage{pgfplots}
\pgfplotsset{compat=1.18}

% Unicode symbols typed directly in prose (e.g. the proof tags (⇒), (⊆))
\usepackage{newunicodechar}
\newunicodechar{⇒}{\ensuremath{\Rightarrow}}
\newunicodechar{⇐}{\ensuremath{\Leftarrow}}
\newunicodechar{⇔}{\ensuremath{\Leftrightarrow}}
\newunicodechar{→}{\ensuremath{\to}}
\newunicodechar{⊆}{\ensuremath{\subseteq}}
\newunicodechar{⊇}{\ensuremath{\supseteq}}
\newunicodechar{∈}{\ensuremath{\in}}
\newunicodechar{≤}{\ensuremath{\le}}
\newunicodechar{≥}{\ensuremath{\ge}}
\newunicodechar{≠}{\ensuremath{\ne}}
\newunicodechar{×}{\ensuremath{\times}}
\newunicodechar{①}{\textcircled{\scriptsize 1}}
\newunicodechar{②}{\textcircled{\scriptsize 2}}
\newunicodechar{③}{\textcircled{\scriptsize 3}}
\newunicodechar{④}{\textcircled{\scriptsize 4}}
\newunicodechar{⑤}{\textcircled{\scriptsize 5}}
\newunicodechar{⑥}{\textcircled{\scriptsize 6}}
\newunicodechar{⑦}{\textcircled{\scriptsize 7}}
\newunicodechar{⑧}{\textcircled{\scriptsize 8}}
\newunicodechar{⑨}{\textcircled{\scriptsize 9}}
\newunicodechar{∗}{\ensuremath{\ast}}
\newunicodechar{−}{\ensuremath{-}}
\newunicodechar{ᵀ}{\texorpdfstring{\ensuremath{^{\top}}}{T}}

% Tables written by Pandoc (column widths use \real from calc)
\usepackage{longtable,booktabs,array,calc}
\providecommand{\tightlist}{\setlength{\itemsep}{0pt}\setlength{\parskip}{0pt}}

% hyperref is loaded by the pandoc template after latex/macros.tex, and it
% reclaims \H (the Hungarian-umlaut accent). Restore the bold matrix H once
% the preamble is over, so \H means the matrix in Chapter 7 onwards.
\AtBeginDocument{\def\H{\mathbf{H}}}

% Let TeX stretch a line rather than run it into the margin. Measured over three
% full builds: 1em removes 61% of the body's overfull boxes (285 -> 110) and
% changes no label number, no page number and no heading break. 3em reaches 38
% but starts re-breaking subsection headings onto two conspicuously loose lines,
% which reads worse than the overrun it fixes. Most of the book's remaining
% overfull boxes are in the Contents, which this does not touch.
\emergencystretch=1em
